% !TeX root = ../thesis.tex
\chapter{Felhasználói dokumentáció}
\label{ch:user}

Az alkalmazás sem telepítést sem regisztrációt nem igényel, teljes egészében böngészőből használható. Minden alkalmazáson belüli tevékenység egy projekthez köthető melyeket a felhasználó kedve szerint hozhat létre, törölhet, módosíthat. Egyszerre több projektet is vezethetünk, a váltás közöttük azonnali.

\section{A Projekt}

Minden projekt \ref{fig:basic-project-structure} három fő elemből áll, ezek a név, a világ és később a szereplők. Új projekt létrehozásakor a projekt nevét és a világának adatait szükséges megadni.

\begin{figure}[h!]
	\centering
	\begin{forest}
		forked edges,
		for tree={edge+={-Latex}},
		[Projekt
			[Név]
			[Világ [Név] [Méret]]
			[Szereplők [\dots]]
		]
	\end{forest}
	\caption{
		Projektstruktúra}
	\label{fig:basic-project-structure}
\end{figure}

\begin{figure}[h!]
	\centering
	\includegraphics[width=0.50\textwidth,scale=1]{create-project-button}
	\caption{
		Új projektet létrehozni a jobb felső sarokban található menü legfelső elemével tehetjük meg}
	\label{fig:create-project-button}
\end{figure}

Új projekt létrehozásakor a projekt készítő dialógust kapjuk melyen lehetőséget kapunk a projektünknek nevet adni és a világbeállításokat megadni, mint a bolygó neve és mérete ami alatt a bolygó sugarát értjük kilóméterben megadva. Valamint egy tetszőleges fekete-fehér magasságtérképet is megadhatunk amivel inicializálhatjuk a bolygó felületét. A vízmagasság alapértelmezetten 40\%-os szürkeségnél van.


\begin{figure}[h!]
	\centering
	\includegraphics[width=0.80\textwidth,scale=1]{create-project-form}
	\caption{
		Új projekt dialógus}
	\label{fig:create-project-form}
\end{figure}


\begin{table}[H]
	\centering
	\begin{tabular}{ | m{0.25\textwidth} | m{0.10\textwidth} | m{0.55\textwidth} | }
		\hline
		\textbf{Mező} & \textbf{Típus} & \textbf{Leírás} \\
		\hline \hline
		\emph{Név} & Szöveg & Kötelező. Egyedi. A projekt azonosítója. \\
		\hline
		\emph{Bolygó név} & Szöveg & Kötelező. Alapértelmezett értéke: "Earth" \\
		\hline
		\emph{Bolygó méret} & Szám & Kötelező. Alapértelmezett értéke: "6371". (A föld sugara kilóméterben)  \\
		\hline
		\emph{Magasságtérkép} & Kép & Opcionális. Egy tetszőleges kezdeti magasságtérkép  \\
		\hline
	\end{tabular}
	\caption{Új projekt sablon}
	\label{tab:create-project-form}
\end{table}

A 'Create' gombra kattintva azonnal megnyílik az elkészített projekt.

\section{Kezelőfelület} \label{section:ui}

Single Page Application
\begin{figure}[h!]
	\centering
	\includegraphics[width=1.0\textwidth,scale=1]{full-screenshot-open-sidebar}
	\caption{
		A kezelőfelület}
	\label{fig:full-screenshot-open-sidebar}
\end{figure}

A kezelőfelület az alábbi elemeket tartalmazza:

\begin{itemize}
	\item Bal oldalt a szereplőtár amibel készíteni és válogatni tudunk a szereplőink között. \ref{section:ui-actors}
	\item Alul az idővonal látható ahol az éppen megtekintett időpontot tudjuk megtekintni/változtatni. Az események idejét szerkeszteni.  \ref{section:ui-timeline}
\end{itemize}

\subsection{Szereplőtár} \label{section:ui-actors}

A szereplőtár legfelső elemével lehet újabb szereplőket hozzáadni a színtérhez. A gombot fogd-és-vidd módon lehet a bolygó valamely pontjára ráhúzni, ami majd a karakter kiindulási pontja lesz. Ez akármikor megváltoztatható.

Ez alatt az elkészített karakterek listája van, mellette az összes karakter számával. Egy karakter itt az alábbi információt tartalmazza:

\begin{itemize}
	\item Név (Akkumulált a jelenlegi időpillanatra)
	\item Szín (Akkumulált a jelenlegi időpillanatra)
	\item Összes történés száma
\end{itemize}

\subsection{Idővonal} \label{section:ui-timeline}

A idővonalat elsősorban három elemre tudjuk bontani:

\subsubsection{A Kurzor}

A vonalzón helyezkedik el. Mindig a jelenlegi időpontra mutat, melyet a kurzoron, fogd-és-vidd módszerrel megváltoztathatunk. Egy fölötte megjelenő buborék a pontos időt mutatja másodperc pontossággal.

\paragraph{A Vonalzó}

Beosztásai napokban értendőek, a kisebb fogai pedig óráknak. A vonalzót fogd és vidd módon tudjuk az időben előre vagy hátra eltolni. Az egérgörgővel pedig az aktív idősávot tudjuk kibővíteni, szűkíteni.

\paragraph{A Karakter Sávok}

Minden karakter saját sávval rendelkezik. Bal oldalt a sávhoz tartozó karakter akkumulált nevét lehet látni. A sávon belül pedig a blokkját. Minden 'blokk' történésekből áll melyeket egy-egy 'cövek' reprezentál.

Egy blokk az alábbi interakciós lehetőségeket nyútja:
\begin{enumerate}
	\item Egy cövekre kattintva a kurzor arra az időpontra ugrik és a cövek kiválasztásra kerül. Ha egy cövek ki van választva mellette egy 'kuka' ikon lehetőséget ad az adott történés eltörlésére. (Megerősítendő)
	\item Aliquam et aliquam turpis, sed fermentum mauris. Nulla eget ex diam.
	\item Donec eget tellus pharetra, semper neque eget, rutrum diam Step~\ref{step:first}.
\end{enumerate}


Donec dapibus sodales ante, at scelerisque nunc laoreet sit amet. Mauris porttitor tincidunt neque, vel ullamcorper neque pulvinar et. Integer eu lorem euismod, faucibus lectus sed, accumsan felis. Nunc ornare mi at augue vulputate, eu venenatis magna mollis. Nunc sed posuere dui, et varius nulla. Sed mollis nibh augue, eget scelerisque eros ornare nec.

\begin{enumerate}
	\item\label{step:first} Donec pretium et quam a cursus. Ut sollicitudin tempus urna et mollis.
	\item Aliquam et aliquam turpis, sed fermentum mauris. Nulla eget ex diam.
	\item Donec eget tellus pharetra, semper neque eget, rutrum diam Step~\ref{step:first}.
\end{enumerate}

Praesent porta, metus eget eleifend consequat, eros ligula eleifend ex, a pellentesque mi est vitae urna. Vivamus turpis nunc, iaculis non leo eget, mattis vulputate tellus. Maecenas rutrum eros sem, pharetra interdum nulla porttitor sit amet. In vitae viverra ante. Maecenas sit amet placerat orci, sed tincidunt velit. Vivamus mattis, enim vel suscipit elementum, quam odio venenatis elit\footnote{Phasellus faucibus varius purus, nec tristique enim porta vitae.}, et mollis nulla nunc a risus. Praesent purus magna, tristique sed lacus sit amet, convallis malesuada magna.

\begin{description}
	\item[Vestibulum venenatis] malesuada enim, ac auctor erat vestibulum et. Phasellus id purus a leo suscipit accumsan.
	\item[Orci varius natoque] penatibus et magnis dis parturient montes, nascetur ridiculus mus. Nullam interdum rhoncus nisl, vel pharetra arcu euismod sagittis. Vestibulum ac turpis auctor, viverra turpis at, tempus tellus.
	\item[Morbi dignissim] erat ut rutrum aliquet. Nulla eu rutrum urna. Integer non urna at mauris scelerisque rutrum sed non turpis.
\end{description}

\subsection{Szoros térközű felsorolások}

Phasellus ultricies, sapien sit amet ultricies placerat, velit purus viverra ligula, id consequat ipsum odio imperdiet enim:
\begin{compactenum}
	\item Maecenas eget lobortis leo.
	\item Donec eget libero enim.
	\item In eu eros a eros lacinia maximus ullamcorper eget augue.
\end{compactenum}

\bigskip

In quis turpis metus. Proin maximus nibh et massa eleifend, a feugiat augue porta. Sed eget est purus. Duis in placerat leo. Donec pharetra eros nec enim convallis:
\begin{compactitem}
	\item Pellentesque odio lacus.
	\item Maximus ut nisl auctor.
	\item Sagittis vulputate lorem.
\end{compactitem}

\bigskip

Vestibulum ante ipsum primis in faucibus orci luctus et ultrices posuere cubilia Curae; Sed lorem libero, dignissim vitae gravida a, ornare vitae est.
\begin{compactdesc}
	\item[Cras maximus] massa commodo pellentesque viverra.
	\item[Morbi sit] amet ante risus. Aliquam nec sollicitudin mauris
	\item[Ut aliquam rhoncus sapien] luctus viverra arcu iaculis posuere
\end{compactdesc}


\section{Képek, ábrák}

Aliquam vehicula luctus mi a pretium. Nulla quam neque, maximus nec velit in, aliquam mollis tortor. Aliquam erat volutpat. Curabitur vitae laoreet turpis. Integer id diam ligula. Nulla sodales purus id mi consequat, eu venenatis odio pharetra. Cras a arcu quam. Suspendisse augue risus, pulvinar a turpis et, commodo aliquet turpis. Nulla aliquam scelerisque mi eget pharetra. Mauris sed posuere elit, ac lobortis metus. Proin lacinia sit amet diam sed auctor. Nam viverra orci id sapien sollicitudin, a aliquam lacus suscipit, Figure~\ref{fig:example-1}:

\begin{figure}[H]
	\centering
	\includegraphics[width=0.6\textwidth,height=100px]{elte_cimer_szines}
	\caption{Quisque ac tincidunt leo}
	\label{fig:example-1}
\end{figure}

\subsection{Képek szegélyezése}

Ut aliquet nec neque eget fermentum. Cras volutpat tellus sed placerat elementum. Quisque neque dui, consectetur nec finibus eget, blandit id purus. Nam eget ipsum non nunc placerat interdum.

\begin{figure}[H]
	\centering
	\includegraphics[width=0.6\textwidth,height=100px,frame]{elte_cimer_szines}
	\caption{Quisque ac tincidunt leo}
\end{figure}

\subsection{Képek csoportosítása}

In non ipsum fermentum urna feugiat rutrum a at odio. Pellentesque habitant morbi tristique senectus et netus et malesuada fames ac turpis egestas. Nulla tincidunt mattis nisl id suscipit. Sed bibendum ac felis sed volutpat. Nam pharetra nisi nec facilisis faucibus. Aenean tristique nec libero non commodo. Nulla egestas laoreet tempus. Nunc eu aliquet nulla, quis vehicula dui. Proin ac risus sodales, gravida nisi vitae, efficitur neque, Figure~\ref{fig:example-2}:

\begin{figure}[H]
	\centering
	\subfigure[Vestibulum quis mattis urna]{
		\includegraphics[width=0.45\linewidth]{elte_cimer_szines}}
	\hspace{5pt}
	\subfigure[Donec hendrerit quis dui sit amet venenatis]{
		\includegraphics[width=0.45\linewidth]{elte_cimer_szines}}
	\caption{Aenean porttitor mi volutpat massa gravida}
	\label{fig:example-2}
\end{figure}

Nam et nunc eget elit tincidunt sollicitudin. Quisque ligula ipsum, tempor vitae tortor ut, commodo rhoncus diam. Pellentesque habitant morbi tristique senectus et netus et malesuada fames ac turpis egestas. Phasellus vehicula quam dui, eu convallis metus porta ac.


\section{Táblázatok}

Nam magna ex, euismod nec interdum sed, sagittis nec leo. Nam blandit massa bibendum mattis tristique. Phasellus tortor ligula, sodales a consectetur vitae, placerat vitae dolor. Aenean consequat in quam ac mollis.

\begin{table}[H]
	\centering
	\begin{tabular}{ | m{0.25\textwidth} | m{0.65\textwidth} | }
		\hline
		\textbf{Phasellus tortor} & \textbf{Aenean consequat} \\
		\hline \hline
		\emph{Sed malesuada} & Aliquam aliquam velit in convallis ultrices. \\
		\hline
		\emph{Purus sagittis} &  Quisque lobortis eros vitae urna lacinia euismod. \\
		\hline
		\emph{Pellentesque} & Curabitur ac lacus pellentesque, eleifend sem ut, placerat enim. Ut auctor tempor odio ut dapibus. \\
		\hline
	\end{tabular}
	\caption{Maecenas tincidunt non justo quis accumsan}
	\label{tab:example-1}
\end{table}

\subsection{Sorok és oszlopok egyesítése}

Mauris a dapibus lectus. Vestibulum commodo nibh ante, ut maximus magna eleifend vel. Integer vehicula elit non lacus lacinia, vitae porttitor dolor ultrices. Vivamus gravida faucibus efficitur. Ut non erat quis arcu vehicula lacinia. Nulla felis mauris, laoreet sed malesuada in, euismod et lacus. Aenean at finibus ipsum. Pellentesque dignissim elit sit amet lacus congue vulputate.

\begin{table}[htb]
	\centering
	\begin{tabular}{ | c | r | r | r | r | r | r | }
		\hline
		\multirow{2}{*}{\textbf{Quisque}} & \multicolumn{2}{ c | }{\textbf{Suspendisse}} & \multicolumn{2}{ c | }{\textbf{Aliquam}} & \multicolumn{2}{ c | }{\textbf{Vivamus}} \\
		\cline{2-7}
		& Proin & Nunc & Proin & Nunc & Proin & Nunc \\
		\hline \hline
		Leo & 2,80 MB & 100\% & 232 KB & 8,09\% & 248 KB & 8,64\% \\
		\hline
		Vel & 9,60 MB & 100\% & 564 KB & 5,74\% & 292 KB & 2,97\% \\
		\hline
		Auge & 78,2 MB & 100\% & 52,3 MB & 66,88\% & 3,22 MB & 4,12\% \\
		\hline
	\end{tabular}
	\caption[Rövid cím a táblázatjegyzékbe]{Vivamus ac arcu fringilla, fermentum neque sed, interdum erat. Mauris bibendum mauris vitae enim mollis, et eleifend turpis aliquet.}
	\label{tab:example-2}
\end{table}

\subsection{Több oldalra átnyúló táblázatok}

Nunc porta placerat leo, sit amet porttitor dui porta molestie. Aliquam at fermentum mi. Maecenas vitae lorem at leo tincidunt volutpat at nec tortor. Vivamus semper lacus eu diam laoreet congue. Vivamus in ipsum risus. Nulla ullamcorper finibus mauris non aliquet. Vivamus elementum rhoncus ex ut porttitor.

\begin{center}
	\begin{longtable}{ | p{0.3\textwidth} | p{0.7\textwidth} | }

		\hline
		\multicolumn{2}{|c|}{\textbf{Praesent aliquam mauris enim}}
		\\ \hline

		\emph{Suspendisse potenti} & \emph{Lorem ipsum dolor sit amet}
		\\ \hline \hline
		\endfirsthead % első oldal fejléce

		\hline
		\emph{Suspendisse potenti} & \emph{Lorem ipsum dolor sit amet}
		\\ \hline \hline
		\endhead % többi oldal fejléce

		\hline
		\endfoot % többi oldal lábléce

		\endlastfoot % utolsó oldal lábléce

		\emph{Praesent}
		& Nulla ultrices et libero sit amet fringilla. Nunc scelerisque ante tempus sapien placerat convallis.
		\\ \hline

		\emph{Luctus}
		& Integer hendrerit erat massa, non hendrerit risus convallis at. Curabitur ultrices, justo in imperdiet condimentum, neque tortor luctus enim, luctus posuere massa erat vitae nibh.
		\\ \hline

		\emph{Egestas}
		& Duis fermentum feugiat augue in blandit. Mauris a tempor felis. Pellentesque ultricies tristique dignissim. Pellentesque aliquam semper tristique. Nam nec egestas dolor. Vestibulum id elit quis enim fringilla tempor eu a mauris. Aliquam vitae lacus tellus. Phasellus mauris lectus, aliquam id leo eget, auctor dapibus magna. Fusce lacinia felis ac elit luctus luctus.
		\\ \hline

		\emph{Dignissim}
		& Praesent aliquam mauris enim, vestibulum posuere massa facilisis in. Suspendisse potenti. Nam quam purus, rutrum eu augue ut, varius vehicula tellus. Fusce dui diam, aliquet sit amet eros at, sollicitudin facilisis quam. Phasellus tempor metus vel augue gravida pretium. Proin aliquam aliquam blandit. Nulla id tempus mi. Fusce in aliquam tortor.
		\\ \hline

		\emph{Pellentesque}
		& Donec felis nibh, imperdiet a arcu non, vehicula gravida nibh. Quisque interdum sapien eu massa commodo, ac elementum felis faucibus.
		\\ \hline

		\emph{Molestie}
		& Cras ullamcorper tellus et auctor ultricies. Maecenas tincidunt euismod lectus nec venenatis. Suspendisse potenti. Pellentesque pretium nunc ut euismod cursus. Nam venenatis condimentum quam. Curabitur suscipit efficitur aliquet. Interdum et malesuada fames ac ante ipsum primis in faucibus.
		\\ \hline

		\emph{Vivamus semper}
		& In purus purus, faucibus eu libero vulputate, tristique sodales nunc. Nulla ut gravida dolor. Fusce vel pellentesque mi, vel efficitur eros. Nunc vitae elit tellus. Sed vestibulum auctor consequat.
		\\ \hline

		\emph{Condimentum}
		& Nulla scelerisque, leo et facilisis pretium, risus enim cursus turpis, eu suscipit ipsum ipsum in mauris. Praesent eget pulvinar ipsum, suscipit interdum nunc. Nam varius massa ut justo ullamcorper sollicitudin. Vivamus facilisis suscipit neque, eu fermentum risus. Ut at mi mauris.
		\\ \hline

		\caption{Praesent ullamcorper consequat tellus ut eleifend}
		\label{tab:example-3}
	\end{longtable}
\end{center}
